\chapterimage{figs/chapterhead.png}
\chapter{Troubleshooting and Known Limitations}
\label{chap:troubleshooting-and-known-limitations}

This chapter documents the current troubleshooting surface of the manual as it
exists in the live repository. The text is grounded in the current GUI shell,
model lifecycle, parser diagnostics, plugin manager, worker boundaries, and
issue-reporting flow. It intentionally stays conservative: when the code gives
only a generic failure message, this chapter says so instead of inventing a
more specific diagnosis.

\noindent\textbf{Objective.} Help readers identify the subsystem behind a
failure, find the current diagnostic entry points, and understand which
limitations are real current-state boundaries rather than temporary editorial
gaps.

\noindent\textbf{Audience.} Users, instructors, operators, and developers who
need a single place to check common failure modes and current support limits.

\noindent\textbf{Scope.} Cover build and startup problems, display issues,
model opening and validation failures, parser and expression errors, plugin
loading issues, worker/tool boundaries, log locations, reporting paths, and
limitations that should be stated explicitly before release.

\noindent\textbf{Status.} Confirmed by code inspection and aligned with the
current repository documentation. The chapter uses structural figure
placeholders until the final screenshots and diagrams are produced.

The chapter follows this sequence:
Figure~\ref{fig:troubleshooting-decision-tree},
Figure~\ref{fig:troubleshooting-error-examples},
Figure~\ref{fig:troubleshooting-diagnostic-flow}, and
Figure~\ref{fig:troubleshooting-log-locations}.

\section{How to triage a failure}
\label{sec:troubleshooting-triage}

The first question is always which subsystem failed:
\begin{itemize}
\item build or configure failure: the toolchain, preset, or dependency set is
usually wrong;
\item startup failure: the runtime environment or display backend is usually
the problem;
\item model open or check failure: the model file, its companion layout file,
or a component validator is usually the cause;
\item expression failure: the parser grammar or an invalid symbol reference is
usually the cause;
\item plugin failure: the plugin file or its current system dependencies are
usually the cause;
\item worker or helper-tool failure: the current scope is often a partially
validated boundary rather than a complete workflow contract.
\end{itemize}

When in doubt, start with the current console, simulation, and report panes,
then inspect the process log described later in this chapter. The manual
favours specific evidence over general reassurance.

\begin{figure}[htbp]
\centering
% FIGURE-SPEC-BEGIN
% ID: fig:troubleshooting-decision-tree
% Source of truth:
%   - ../README.md
%   - ../docs/ai_assistants/STATUS.md
%   - ../docs/ai_assistants/runbooks/LOCAL_AGENT.md
%   - ../source/applications/gui/genesys/main.cpp
%   - ../source/applications/gui/genesys/controllers/ModelLifecycleController.cpp
%   - ../source/applications/shell/GenesysShell.cpp
% Purpose:
%   Show a conservative decision tree that routes the reader from symptom class to the current diagnostic surface.
% Required visual content:
%   - build/configure branch
%   - startup/display branch
%   - model open/check branch
%   - parser branch
%   - plugin branch
%   - worker/helper-tool branch
% Accessibility description:
%   A top-down flowchart that helps the reader pick the next diagnostic step without guessing.
% Validation:
%   The branches must reflect the current code and documentation boundaries rather than future support plans.
% FIGURE-SPEC-END
\figureplaceholder{figs/generated/troubleshooting-decision-tree}
\caption{Current troubleshooting decision tree.}
\label{fig:troubleshooting-decision-tree}
\end{figure}

Figure~\ref{fig:troubleshooting-decision-tree} should be the quickest way to
route a new problem to the right subsystem.

\section{Build and startup problems}
\label{sec:troubleshooting-build-startup}

\subsection{Toolchain or preset mismatch}
\label{subsec:troubleshooting-toolchain}

\begin{description}
\item[Symptom.] Configuration or build stops before the application starts.
\item[Likely cause.] The current CMake preset does not match the available
toolchain, Qt6 setup, or build directory state.
\item[Diagnosis.] Re-run the exact preset from the repository root and compare
the reported compiler, CMake, and Qt versions against the current baseline in
the repository README and status documents.
\item[Recovery.] Use the canonical preset for the target you are trying to
build, then rebuild from a clean configuration directory if the previous cache
was produced with a different toolchain.
\item[Current limitation.] A successful configure does not prove that the
runtime workflow is healthy; it only proves the selected preset configured.
\end{description}

\subsection{No display or invalid display backend}
\label{subsec:troubleshooting-display}

\begin{description}
\item[Symptom.] The GUI does not appear, exits immediately, or reports a Qt
display connection failure on a headless host.
\item[Likely cause.] There is no usable display server for the current session.
\item[Diagnosis.] Check whether the session has an interactive display. For
headless validation, use a private Xvfb or equivalent virtual desktop and keep
TCP listening disabled.
\item[Recovery.] Relaunch the GUI in a valid graphical session, or reproduce
the launch under the same virtual-desktop setup used by the local runbook.
\item[Current limitation.] Startup under Xvfb is only startup evidence; it is
not proof of end-to-end interaction or scientific correctness.
\end{description}

The visual summary of these two failure classes appears in
Figure~\ref{fig:troubleshooting-error-examples}.

\begin{figure}[htbp]
\centering
% FIGURE-SPEC-BEGIN
% ID: fig:troubleshooting-error-examples
% Source of truth:
%   - ../source/applications/gui/genesys/main.cpp
%   - ../source/applications/gui/genesys/controllers/ModelLifecycleController.cpp
%   - ../source/applications/gui/genesys/mainwindow.cpp
%   - ../source/applications/shell/GenesysShell.cpp
% Purpose:
%   Show the most common user-visible error messages for startup, model opening, model checking, and plugin auto-load problems.
% Required visual content:
%   - a build/startup failure example
%   - the generic model-open warning
%   - the generic model-check failure dialog
%   - the shell plugin auto-load warning
% Accessibility description:
%   A montage of UI and terminal messages that match current code paths.
% Validation:
%   The text shown must match messages actually emitted by the current code paths.
% FIGURE-SPEC-END
\figureplaceholder{figs/generated/troubleshooting-error-examples}
\caption{Representative current error messages.}
\label{fig:troubleshooting-error-examples}
\end{figure}

\section{Model and parser problems}
\label{sec:troubleshooting-model-parser}

\subsection{A model does not open}
\label{subsec:troubleshooting-open-model}

\begin{description}
\item[Symptom.] Opening a model ends with the generic warning that the model
could not be opened.
\item[Likely cause.] The selected file is unreadable, not a valid model file,
or does not match the companion graphical layout expected by the open flow.
\item[Diagnosis.] The current open flow prefers a companion \texttt{.gui}
file when a \texttt{.gen} file is selected. If the open action fails, inspect
the exact file that was chosen and check whether the matching layout file is
present.
\item[Recovery.] Reopen the companion \texttt{.gui} when it exists, or reopen a
known-good model from the \texttt{models/} tree and confirm that the current
file path is valid.
\item[Current limitation.] The dialog message is intentionally generic; the
actual cause has to be recovered from the file choice and the console trace.
\end{description}

\subsection{Model check fails}
\label{subsec:troubleshooting-model-check}

\begin{description}
\item[Symptom.] The model check dialog reports failure instead of the success
message.
\item[Likely cause.] A component validator rejected a parameter, a connection,
or a model rule that is checked during the current model validation pass.
\item[Diagnosis.] Inspect the console. The current dialog only tells the user
that the model has errors and to see the console for more information.
\item[Recovery.] Fix the component or property that produced the console
message, then run model check again.
\item[Current limitation.] Model check success only proves the current
validation pass succeeded; it does not prove that the model is scientifically
validated or that every future scenario is safe.
\end{description}

\subsection{Parser or expression errors}
\label{subsec:troubleshooting-parser}

\begin{description}
\item[Symptom.] Expression evaluation in the parser checker or a property
editor fails and the console shows an error line instead of a numeric result.
\item[Likely cause.] The expression does not match the current grammar, a
symbol is missing from the current model context, or the parser returned a
syntax error.
\item[Diagnosis.] Use the parser grammar checker dialog and the console output.
The checker currently reports either the real parser message or a fallback
\texttt{Unknown parser error.} message when no more specific detail is
available.
\item[Recovery.] Rewrite the expression to match the grammar currently
implemented in the parser, then evaluate it again in the same model context.
\item[Current limitation.] The parser diagnostics are helpful but not a full
formal proof of semantic correctness for a model.
\end{description}

\section{Plugins, worker, and helper tools}
\label{sec:troubleshooting-plugins-worker-tools}

\subsection{Plugin load or dependency problems}
\label{subsec:troubleshooting-plugins}

\begin{description}
\item[Symptom.] The plugin manager shows blocked plugins, or the shell reports
that auto-loading plugins from \texttt{autoloadplugins.txt} failed.
\item[Likely cause.] The plugin file is missing, incomplete, structurally
invalid, or still depends on system packages that are not installed.
\item[Diagnosis.] Use the plugin manager's loaded and problem tabs, then check
the current dependency preflight and the system package state. The shell
starts even when the auto-load list is not usable, so a plugin warning alone
does not imply a fatal application failure.
\item[Recovery.] Load the plugin list that matches the current checkout, or
resolve the blocked plugin dependencies through the plugin manager if the
system package prerequisites are available.
\item[Current limitation.] Plugin loading still follows the current static
build graph and the existing deployment/search contract; do not assume a future
dynamic-plugin boundary or a fully automatic recovery path.
\end{description}

\subsection{Worker reachability and scope}
\label{subsec:troubleshooting-worker}

\begin{description}
\item[Symptom.] The worker does not answer, or a worker-based workflow cannot
complete from outside the local machine.
\item[Likely cause.] The worker is not running, is bound outside the expected
local test setup, or is being used outside the trusted boundary currently
documented by the repository.
\item[Diagnosis.] Confirm the worker target and launch path first, then inspect
its local reachability from the same machine or private network segment used
for development.
\item[Recovery.] Use the current local/private deployment pattern only. Do not
try to turn this into a public service as a troubleshooting shortcut.
\item[Current limitation.] The worker is not approved for direct public
Internet exposure, and its full security profile is not yet the supported
production profile.
\end{description}

\subsection{Standalone tools that start but do not complete a workflow}
\label{subsec:troubleshooting-tools}

\begin{description}
\item[Symptom.] A helper window opens, but the intended workflow does not look
complete or reproducible yet.
\item[Likely cause.] The tool is only startup-validated, or the specific
analysis/optimization/AI workflow is still outside the fully validated scope.
\item[Diagnosis.] Distinguish startup from workflow validation. Check whether
the current tool is one of the standalone helper front ends or a partially
validated workflow behind them.
\item[Recovery.] Reduce the problem to the smallest current supported path and
capture the exact inputs, outputs, and logs that were produced.
\item[Current limitation.] Startup evidence does not prove complete functional
or scientific correctness for the helper tools.
\end{description}

Figure~\ref{fig:troubleshooting-diagnostic-flow} summarizes the current
diagnostic path across these subsystems.

\begin{figure}[htbp]
\centering
% FIGURE-SPEC-BEGIN
% ID: fig:troubleshooting-diagnostic-flow
% Source of truth:
%   - ../source/applications/gui/genesys/controllers/DialogUtilityController.cpp
%   - ../source/applications/gui/genesys/controllers/ModelLifecycleController.cpp
%   - ../source/applications/gui/genesys/dialogs/dialogpluginmanager.cpp
%   - ../source/applications/gui/genesys/mainwindow.cpp
%   - ../source/applications/shell/GenesysShell.cpp
% Purpose:
%   Show how the reader moves from symptom to console, dialog, or subsystem-specific recovery path.
% Required visual content:
%   - model open/check branch
%   - parser branch
%   - plugin branch
%   - worker/tool branch
%   - log/reporting branch
% Accessibility description:
%   A flowchart that turns the chapter's prose into a concrete debugging sequence.
% Validation:
%   The flow must not claim more validation than the current code and status documents support.
% FIGURE-SPEC-END
\figureplaceholder{figs/generated/troubleshooting-diagnostic-flow}
\caption{Subsystem-oriented diagnostic flow.}
\label{fig:troubleshooting-diagnostic-flow}
\end{figure}

\section{Logs and reporting}
\label{sec:troubleshooting-logs}

The current GUI application installs a log file in the working directory named
\texttt{crashesAndLogs.log}. The application appends Qt message output and
startup diagnostics to that file as soon as the GUI process starts.

The issue-report workflow in the About menu can attach the tail of that log
file together with the current console, simulation, reports, and optionally the
current model text. That makes \texttt{crashesAndLogs.log} the first file to
check when a GUI problem needs to be reported or reproduced.

For terminal workflows, capture the exact command and the full terminal output
instead of relying on memory. For model problems, also capture the model file
that was opened and whether the companion \texttt{.gui} file was present.

\begin{figure}[htbp]
\centering
% FIGURE-SPEC-BEGIN
% ID: fig:troubleshooting-log-locations
% Source of truth:
%   - ../source/applications/gui/genesys/main.cpp
%   - ../source/applications/gui/genesys/controllers/DialogUtilityController.cpp
%   - ../source/applications/gui/genesys/dialogs/DialogReportIssue.cpp
% Purpose:
%   Show where the current GUI writes its process log and what the issue-report path can attach.
% Required visual content:
%   - the working-directory log file
%   - the About/Report Issue entry point
%   - the attach-log-tail option
%   - the console/simulation/reports panes
% Accessibility description:
%   A compact diagram that points to the current diagnostic evidence locations.
% Validation:
%   The location and attachments must match the current GUI startup and reporting code.
% FIGURE-SPEC-END
\figureplaceholder{figs/generated/troubleshooting-log-locations}
\caption{Current log and reporting locations.}
\label{fig:troubleshooting-log-locations}
\end{figure}

\section{Known limitations}
\label{sec:troubleshooting-known-limitations}

The following boundaries are current state, not promises of future support:
\begin{itemize}
\item startup does not prove complete workflow correctness;
\item helper-tool startup does not prove the associated analysis or
optimization workflow is scientifically validated;
\item model validation failures are reported through generic dialogs with
console details, not through a fully structured problem taxonomy;
\item the shell auto-load contract around \texttt{autoloadplugins.txt} remains a
current deployment and fallback gap;
\item the worker must remain in the current local or private-network boundary;
\item numerical, statistical, and biological results must still be treated as
bounded evidence unless they are explicitly validated elsewhere;
\item a reported problem should be read against the current code and status,
not as a guarantee of future support policy.
\end{itemize}

This chapter exists so the manual can say, clearly and conservatively, where
the current repository stops today.
